\begin{ejer}
Demuestra que $V(I)=V(\sqrt{I})$ y, más aún, que
\[
V(I)\subseteq V(J)\quad\Longleftrightarrow\quad \sqrt{J}\subseteq\sqrt{I}.
\]
\end{ejer}
\begin{sol}
Antes que nada, recordemos que
\[
\sqrt{I}=\bigcap_{\mathfrak{p}\in V(I)}\mathfrak{p}.
\]
\begin{itemize}
    \item Por mera definición de intersección:
    \[
    \mathfrak{p}_0\en V(I)
    \quad\Longleftrightarrow\quad
    \mathfrak{p}_0\supseteq\bigcap_{\mathfrak{p}\in V(I)}\mathfrak{p}=\sqrt{I}
    \quad\Longleftrightarrow\quad
    \mathfrak{p}_0\en V(\sqrt{I}).
    \]

    \item Si $\sqrt{J}\subseteq\sqrt{I}$, se tiene $V(\sqrt{I})\subseteq V(\sqrt{J})$; gracias a lo recién probado, esto significa $V(I)\subseteq V(J)$. La otra implicación es una simple propiedad de la intersección:
    \[
    V(I)\subseteq V(J)
    \quad\Longrightarrow\quad
    \sqrt{J}=\bigcap_{\mathfrak{p}\in V(J)}\mathfrak{p}
    \subseteq\bigcap_{\mathfrak{p}\in V(I)}\mathfrak{p}=\sqrt{I}.
    \]
\end{itemize}
\end{sol}